\begin{textsample}{POS Dim 1 – 2014 – Score 21.00 – 2014/t001874.txt}  \label{ex:f1_pos_046}
What ' s \textbf{the} scariest thing you ' ve ever done?
Or another way to say it is, what ' s \textbf{the} most dangerous thing that you ' ve ever done?
And why did you do it?
I know what \textbf{the} most dangerous thing is that I ' ve ever done because NASA does \textbf{the} math.
You look back to \textbf{the} first five shuttle launches, \textbf{the} odds of a catastrophic event during \textbf{the} first five shuttle launches was one in nine.
And even when I first flew in \textbf{the} shuttle back in 1995, 74 shuttle flight, \textbf{the} odds were still now that we look back about one in 38 or so — one in 35, one in 40.
Not great odds, so it ' s a really interesting day when you wake up at \textbf{the} Kennedy Space Center and you ' re going to go to space that day because you realize by \textbf{the} end of \textbf{the} day you ' re either going to be floating effortlessly, gloriously in space, or you ' ll be dead.
You go into, at \textbf{the} Kennedy Space Center, \textbf{the} suit-up room, \textbf{the} same room that our childhood heroes got dressed in, that Neil Armstrong and Buzz Aldrin got suited in to go ride \textbf{the} Apollo rocket to \textbf{the} moon.
And I got my pressure suit built around me and rode down outside in \textbf{the} van heading out to \textbf{the} launchpad — in \textbf{the} Astro van — heading out to \textbf{the} launchpad, and as you come around \textbf{the} corner at \textbf{the} Kennedy Space Center, it ' s normally predawn, and in \textbf{the} distance, lit up by \textbf{the} huge xenon lights, is your spaceship — \textbf{the} vehicle that is going to take you off \textbf{the} \textbf{planet}. \textbf{The} crew is sitting in \textbf{the} Astro van sort of hushed, almost holding hands, looking at that as it gets bigger and bigger.
We ride \textbf{the} elevator up and we crawl in, on your hands and knees into \textbf{the} spaceship, one at a time, and you worm your way up into your chair and plunk yourself down on your back.
And \textbf{the} hatch is closed, and suddenly, what has been a lifetime of both dreams and denial is becoming real, something that I dreamed about, in fact, that I chose to do when I was nine years old, is now suddenly within not too many minutes of actually happening.
In \textbf{the} astronaut business — \textbf{the} shuttle is a very complicated vehicle ; it ' s \textbf{the} most complicated flying machine ever built.
And in \textbf{the} astronaut business, we have a saying, which is, there is no problem so bad that you can ' t make it worse.
And so you ' re very conscious in \textbf{the} cockpit ; you ' re thinking about all of \textbf{the} things that you might have to do, all \textbf{the} switches and all \textbf{the} wickets you have to go through.
And as \textbf{the} time gets closer and closer, this excitement is building.
And then about three and a half minutes before launch, \textbf{the} huge nozzles on \textbf{the} back, like \textbf{the} size of big church bells, swing back and forth and \textbf{the} mass of them is such that it sways \textbf{the} whole vehicle, like \textbf{the} vehicle is alive underneath you, like an \textbf{elephant} getting up off its knees or something.
And then about 30 seconds before launch, \textbf{the} vehicle is completely alive — it is ready to go — \textbf{the} APUs are running, \textbf{the} \textbf{computers} are all self-contained, it ' s ready to leave \textbf{the} \textbf{planet}.
And 15 seconds before launch, this happens : Voice : 12, 11, 10, nine, eight, seven, six — — start, two, one, booster ignition, and liftoff of \textbf{the} space shuttle Discovery, returning to \textbf{the} space station, paving \textbf{the} way...
Chris Hadfield : It is incredibly powerful to be on board one of these things.
You are in \textbf{the} grip of something that is vastly more powerful than yourself.
It ' s shaking you so hard you can ' t focus on \textbf{the} instruments in front of you.
It ' s like you ' re in \textbf{the} jaws of some enormous dog and there ' s a foot in \textbf{the} small of your back pushing you into space, accelerating wildly straight up, shouldering your way through \textbf{the} air, and you ' re in a very complex place — paying attention, watching \textbf{the} vehicle go through each one of its wickets with a steadily increasing smile on your face.
After two minutes, those solid rockets explode off and then you just have \textbf{the} liquid engines, \textbf{the} hydrogen and oxygen, and it ' s as if you ' re in a dragster with your foot to \textbf{the} floor and accelerating like you ' ve never accelerated.
You get lighter and lighter, \textbf{the} force gets on us heavier and heavier.
It feels like someone ' s pouring cement on you or something.
Until finally, after about eight minutes and 40 seconds or so, we are finally at exactly \textbf{the} right \textbf{altitude}, exactly \textbf{the} right speed, \textbf{the} right direction, \textbf{the} engine shut off, and we ' re weightless.
And we ' re alive.
It ' s an amazing experience.
But why would we take that risk?
Why would you do something that dangerous?
In my case \textbf{the} answer is fairly straightforward.
I was inspired as a youngster that this was what I wanted to do.
I watched \textbf{the} first people walk on \textbf{the} moon and to me, it was just an obvious thing — I want to somehow turn myself into that.
But \textbf{the} real question is, how do you deal with \textbf{the} danger of it and \textbf{the} fear that comes from it?
How do you deal with fear versus danger?
And having \textbf{the} goal in mind, thinking about where it might lead, directed me to a life of looking at all of \textbf{the} small details to allow this to become possible, to be able to launch and go help build a space station where you are on board a million-pound creation that ' s going around \textbf{the} world at five miles a second, eight kilometers a second, around \textbf{the} world 16 times a day, with experiments on board that are teaching us what \textbf{the} substance of \textbf{the} \textbf{universe} is made of and running 200 experiments inside.
But maybe even more importantly, allowing us to see \textbf{the} world in a way that is impossible through any other means, to be able to look down and have — if your jaw could drop, it would — \textbf{the} jaw-dropping gorgeousness of \textbf{the} turning orb like a self-propelled art gallery of \textbf{fantastic}, constantly changing beauty that is \textbf{the} world itself.
And you see, because of \textbf{the} speed, a sunrise or a sunset every 45 minutes for half a year.
And \textbf{the} most magnificent part of all that is to go outside on a spacewalk.
You are in a one- person spaceship that is your spacesuit, and you ' re going through space with \textbf{the} world.
It ' s an entirely different perspective, you ' re not looking up at \textbf{the} \textbf{universe}, you and \textbf{the} Earth are going through \textbf{the} \textbf{universe} together.
And you ' re holding on with one hand, looking at \textbf{the} world turn beside you.
It ' s roaring silently with color and texture as it pours by mesmerizingly next to you.
And if you can tear your eyes away from that and you look under your arm down at \textbf{the} rest of everything, it ' s unfathomable blackness, with a texture you feel like you could stick your hand into. and you are holding on with one hand, one link to \textbf{the} other seven billion people.
And I was outside on my first spacewalk when my \textbf{left} eye went blind, and I didn ' t know why.
Suddenly my \textbf{left} eye slammed shut in great pain and I couldn ' t figure out why my eye wasn ' t working.
I was thinking, what do I do next?
I thought, well maybe that ' s why we have two eyes, so I kept working.
But unfortunately, without gravity, tears don ' t fall.
So you just get a bigger and bigger ball of whatever that is mixed with your tears on your eye until eventually, \textbf{the} ball becomes so big that \textbf{the} \textbf{surface} tension takes it across \textbf{the} bridge of your nose like a tiny little waterfall and goes"goosh"into your other eye, and now I was completely blind outside \textbf{the} spaceship.
So what ' s \textbf{the} scariest thing you ' ve ever done?
Maybe it ' s \textbf{spiders}.
A lot of people are afraid of \textbf{spiders}.
I think you should be afraid of \textbf{spiders} — \textbf{spiders} are creepy and they ' ve got long, hairy legs, and \textbf{spiders} like this one, \textbf{the} brown recluse — it ' s horrible.
If a brown recluse \textbf{bites} you, you end with one of these horrible, big necrotic things on your leg and there might be one right now sitting on \textbf{the} chair behind you, in fact.
And how do you know?
And so a \textbf{spider} lands on you, and you go through this great, spasmy attack because \textbf{spiders} are scary.
But then you could say, well is there a brown recluse sitting on \textbf{the} chair beside me or not?
I don ' t know.
Are there brown recluses here?
So if you actually do \textbf{the} research, you find out that in \textbf{the} world there are about 50, 000 different types of \textbf{spiders}, and there are about two dozen that are venomous out of 50, 000.
And if you ' re in Canada, because of \textbf{the} cold winters here in B.
C., there ' s about 720, 730 different types of \textbf{spiders} and there ' s one — one — that is venomous, and its venom isn ' t even fatal, it ' s just kind of like a nasty sting.
And that \textbf{spider} — not only that, but that \textbf{spider} has beautiful markings on it, it ' s like"I ' m dangerous.
I got a big radiation symbol on my back, it ' s \textbf{the} black widow."So, if you ' re even slightly careful you can avoid running into \textbf{the} one \textbf{spider} — and it lives close \textbf{the} ground, you ' re walking along, you are never going to go through a \textbf{spider} \textbf{web} where a black widow \textbf{bites} you. \textbf{Spider} \textbf{webs} like this, it doesn ' t build those, it builds them down in \textbf{the} corners.
And its a black widow because \textbf{the} female \textbf{spider} eats \textbf{the} male ; it doesn ' t care about you.
So in fact, \textbf{the} next time you walk into a spiderweb, you don ' t need to panic and go with your caveman reaction. \textbf{The} danger is entirely different than \textbf{the} fear.
How do you get around it, though?
How do you change your behavior?
Well, next time you see a spiderweb, have a good look, make sure it ' s not a black widow \textbf{spider}, and then walk into it.
And then you see another spiderweb and walk into that one.
It ' s just a little bit of fluffy stuff.
It ' s not a big deal.
And \textbf{the} \textbf{spider} that may come out is no more threat to you than a lady bug or a butterfly.
And then I guarantee you if you walk through 100 spiderwebs you will have changed your fundamental human behavior, your caveman reaction, and you will now be able to walk in \textbf{the} park in \textbf{the} morning and not worry about that spiderweb — or into your grandma ' s attic or whatever, into your own basement.
And you can apply this to anything.
If you ' re outside on a spacewalk and you ' re blinded, your natural reaction would be to panic, I think.
It would make you nervous and worried.
But we had considered all \textbf{the} venom, and we had practiced with a whole variety of different spiderwebs.
We knew everything there is to know about \textbf{the} spacesuit and we trained \textbf{underwater} thousands of times.
And we don ' t just practice things going right, we practice things going wrong all \textbf{the} time, so that you are constantly walking through those spiderwebs.
And not just \textbf{underwater}, but also in virtual reality labs with \textbf{the} helmet and \textbf{the} gloves so you feel like it ' s realistic.
So when you finally actually get outside on a spacewalk, it feels much different than it would if you just went out first time.
And even if you ' re blinded, your natural, panicky reaction doesn ' t happen.
Instead you kind of look around and go,"Okay, I can ' t see, but I can hear, I can talk, Scott Parazynski is out here with me.
He could come over and help me."We actually practiced incapacitated crew rescue, so he could float me like a blimp and stuff me into \textbf{the} airlock if he had to.
I could find my own way back.
It ' s not nearly as big a deal.
And actually, if you keep on crying for a while, whatever that gunk was that ' s in your eye starts to dilute and you can start to see again, and Houston, if you negotiate with them, they will let you then keep working.
We finished everything on \textbf{the} spacewalk and when we came back inside, Jeff got some cotton batting and took \textbf{the} crusty stuff around my eyes, and it turned out it was just \textbf{the} anti-fog, sort of a mixture of \textbf{oil} and soap, that got in my eye.
And now we use Johnson ' s No More Tears, which we probably should ' ve been using right from \textbf{the} very beginning.
But \textbf{the} key to that is by looking at \textbf{the} difference between perceived danger and actual danger, where is \textbf{the} real risk?
What is \textbf{the} real thing that you should be afraid of?
Not just a generic fear of bad things happening.
You can fundamentally change your reaction to things so that it allows you to go places and see things and do things that otherwise would be completely denied to you... where you could see \textbf{the} hardpan south of \textbf{the} Sahara, or you can see New York City in a way that is almost dreamlike, or \textbf{the} unconscious gingham of Eastern Europe fields or \textbf{the} Great Lakes as a \textbf{collection} of small puddles.
You can see \textbf{the} fault lines of San Francisco and \textbf{the} way \textbf{the} water pours out under \textbf{the} bridge, just entirely different than any other way that you could have if you had not found a way to conquer your fear.
You see a beauty that otherwise never would have happened.
It ' s time to come home at \textbf{the} end.
This is our spaceship, \textbf{the} Soyuz, that little one.
Three of us climb in, and then this spaceship detaches from \textbf{the} station and falls into \textbf{the} atmosphere.
These two parts here actually melt, we jettison them and they burn up in \textbf{the} atmosphere. \textbf{The} only part that survives is \textbf{the} little \textbf{bullet} that we ' re riding in, and it falls into \textbf{the} atmosphere, and in essence you are riding a meteorite home, and riding meteorites is scary, and it ought to be.
But instead of riding into \textbf{the} atmosphere just screaming, like you would if suddenly you found yourself riding a meteorite back to Earth — — instead, 20 years previously we had started studying Russian, and then once you learn Russian, then we learned orbital mechanics in Russian, and then we learned vehicle control theory, and then we got into \textbf{the} simulator and practiced over and over and over again.
And in fact, you can fly this meteorite and steer it and land in about a 15-kilometer circle anywhere on \textbf{the} Earth.
So in fact, when our crew was coming back into \textbf{the} atmosphere inside \textbf{the} Soyuz, we weren ' t screaming, we were laughing ; it was fun.
And when \textbf{the} great big parachute opened, we knew that if it didn ' t open there ' s a second parachute, and it runs on a nice little clockwork mechanism.
So we came back, we came thundering back to Earth and this is what it looked like to land in a Soyuz, in Kazakhstan.
Reporter : And you can see one of those \textbf{search} and recovery helicopters, once again that helicopter part of dozen such Russian Mi- 8 helicopters.
Touchdown — 3:14 and 48 seconds, a. m.
Central Time.
CH : And you roll to a stop as if someone threw your spaceship at \textbf{the} ground and it tumbles end over end, but you ' re ready for it you ' re in a custom-built seat, you know how \textbf{the} shock absorber works.
And then eventually \textbf{the} Russians reach in, drag you out, plunk you into a chair, and you can now look back at what was an incredible experience.
You have taken \textbf{the} dreams of that nine-year-old boy, which were impossible and dauntingly scary, dauntingly \textbf{terrifying}, and put them into practice, and figured out a way to reprogram yourself, to change your primal fear so that it allowed you to come back with a set of experiences and a level of inspiration for other people that never could have been possible otherwise.
Just to finish, they asked me to play that guitar.
I know this song, and it ' s really a tribute to \textbf{the} genius of David Bowie himself, but it ' s also, I think, a reflection of \textbf{the} fact that we are not machines exploring \textbf{the} \textbf{universe}, we are people, and we ' re taking that ability to adapt and that ability to understand and \textbf{the} ability to take our own self-perception into a new place. — This is Major Tom to ground control — — I ' ve left forevermore — — And I ' m floating in a most \textbf{peculiar} way — — And \textbf{the} stars look very different today — — For here am I floating in \textbf{the} tin can — — A last glimpse of \textbf{the} world — — Planet Earth is \textbf{blue} and there ' s so much left to do — Fear not.
That ' s very nice of you.
Thank you very much.
Thank you.

% matched lemmas: altitude, bite, blue, bullet, collection, computer, elephant, fantastic, left, oil, peculiar, planet, search, spider, surface, terrifying, the, underwater, universe, web
\end{textsample}
